%========================================
% MatinBook - Stage 06: Theorem System Test
% Test: All Theorem Environments + Boxes + References
% Purpose: Validate theorem, definition, example, proof, and all box types
%========================================

% Add parent directory to search path
\makeatletter
\def\input@path{{../}}
\makeatother

\documentclass{matinbook}

\title{آزمون سیستم قضایا و محیط‌های علمی}
\author{تیم توسعه MatinBook}
\date{تابستان ۱۴۰۵}

\begin{document}
    
    %----------------------------------------
    % Title Page
    %----------------------------------------
    \maketitle
    
    %----------------------------------------
    % Chapter 1: Number Theory
    %----------------------------------------
    \chapter{نظریه اعداد}
    
    \section{اعداد اول}
    
    \begin{definition}[عدد اول]
        عدد طبیعی $n > 1$ را \textbf{اول}\index{عدد اول} گوییم
        هرگاه تنها مقسوم‌علیه‌های مثبت آن ۱ و خود $n$ باشند.
        اعدادی که اول نباشند را \textbf{مرکب} گوییم.
        \label{def:prime}
    \end{definition}
    
    طبق \defref{def:prime}، عدد ۱ نه اول است و نه مرکب.
    این قرارداد برای سادگی در بیان قضایا اتخاذ شده است.
    
    \begin{example}
        اعداد اول کوچکتر از ۵۰ عبارتند از:
        \[
        2, 3, 5, 7, 11, 13, 17, 19, 23, 29, 31, 37, 41, 43, 47
        \]
        \label{ex:primes-under-50}
    \end{example}
    
    \begin{example}
        عدد ۶ اول نیست (مرکب است) زیرا $6 = 2 \times 3$.
        مقسوم‌علیه‌های عدد ۶ عبارتند از: ۱، ۲، ۳، ۶.
        \label{ex:not-prime}
    \end{example}
    
    \section{قضیه اساسی حساب}
    
    \begin{theorem}[قضیه اساسی حساب]
        هر عدد طبیعی $n > 1$ را می‌توان به‌طور یکتا
        (صرف‌نظر از ترتیب عوامل) به صورت حاصلضرب اعداد اول نوشت:
        \[
        n = p_{1}^{e_{1}} \cdot p_{2}^{e_{2}} \cdots p_{k}^{e_{k}}
        \]
        که در آن $p_{1} < p_{2} < \cdots < p_{k}$ اعداد اول
        و $e_{i} \geq 1$ اعداد صحیح هستند.
        \label{thm:fundamental}
    \end{theorem}
    
    طبق \thmref{thm:fundamental}، هر عدد طبیعی بزرگتر از ۱
    دارای یک تجزیه یکتا به عوامل اول است. برای مثال:
    \[
    60 = 2^{2} \times 3 \times 5
    \]
    
    \begin{proof}
        ابتدا \textbf{وجود} تجزیه را با استقرای قوی اثبات می‌کنیم.
        
        \textbf{پایه:} عدد ۲ اول است و تجزیه آن خودش می‌باشد.
        
        \textbf{گام استقرا:} فرض کنید هر عدد $m$ با $2 \leq m < n$
        قابل تجزیه به عوامل اول باشد. برای $n$ دو حالت داریم:
        \begin{enumerate}
            \item اگر $n$ اول باشد، تجزیه آن خودش است.
            \item اگر $n$ مرکب باشد، $n = a \cdot b$ که $2 \leq a, b < n$.
            طبق فرض استقرا، $a$ و $b$ هر یک به عوامل اول تجزیه می‌شوند.
            از ضرب این دو تجزیه، تجزیه $n$ به‌دست می‌آید.
        \end{enumerate}
        
        برای اثبات \textbf{یکتایی} از لم اقلیدس (\lemref{lem:euclid})
        استفاده می‌شود که در ادامه خواهیم دید.
    \end{proof}
    
    \section{لم اقلیدس}
    
    \begin{lemma}[لم اقلیدس]
        اگر $p$ عددی اول باشد و $p \mid ab$، آنگاه
        $p \mid a$ یا $p \mid b$.
        \label{lem:euclid}
    \end{lemma}
    
    \begin{proof}
        فرض کنید $p \nmid a$. آنگاه $\gcd(p, a) = 1$
        (چون $p$ اول است). طبق قضیه بزو (\lr{Bezout})،
        اعداد صحیح $x$ و $y$ وجود دارند که:
        \[
        px + ay = 1
        \]
        با ضرب طرفین در $b$:
        \[
        pbx + aby = b
        \]
        از آنجا که $p \mid p$ و $p \mid ab$،
        نتیجه می‌شود $p \mid b$.
    \end{proof}
    
    \begin{corollary}
        اگر $p$ عددی اول باشد و $p \mid a_{1}a_{2} \cdots a_{n}$،
        آنگاه $p$ حداقل یکی از $a_{i}$ها را می‌شمارد
        (بر آن بخش‌پذیر است).
        \label{cor:euclid-general}
    \end{corollary}
    
    \begin{proposition}
        کوچکترین مقسوم‌علیه بزرگتر از ۱ هر عدد طبیعی $n > 1$،
        یک عدد اول است.
        \label{prop:smallest-divisor}
    \end{proposition}
    
    \section{نکات مهم}
    
    \begin{remark}
        لم اقلیدس (\lemref{lem:euclid}) یکی از پایه‌های
        اساسی نظریه اعداد است و در اثبات بسیاری از قضایا
        از جمله یکتایی تجزیه در قضیه اساسی حساب
        (\thmref{thm:fundamental}) به کار می‌رود.
        \label{rem:euclid-importance}
    \end{remark}
    
    \begin{remark}
        عدد ۲ تنها عدد اول زوج است. تمام اعداد اول دیگر فرد هستند،
        زیرا اگر یک عدد زوج بزرگتر از ۲ اول می‌بود،
        باید بر ۲ بخش‌پذیر می‌بود که با تعریف عدد اول
        در تناقض است.
        \label{rem:two-is-special}
    \end{remark}
    
    %----------------------------------------
    % Chapter 2: Group Theory
    %----------------------------------------
    \chapter{مبانی نظریه گروه‌ها}
    
    \section{تعریف گروه}
    
    \begin{definition}[گروه]
        یک \textbf{گروه}\index{گروه} $(G, *)$ شامل یک مجموعه
        ناتهی $G$ و یک عمل دوتایی $*$ روی $G$ است
        که در چهار شرط زیر صدق کند:
        \begin{enumerate}
            \item \textbf{بسته بودن:}
            $\forall a, b \in G: a * b \in G$
            \item \textbf{شرکت‌پذیری:}
            $\forall a, b, c \in G: (a * b) * c = a * (b * c)$
            \item \textbf{وجود عضو خنثی:}
            $\exists e \in G: \forall a \in G: e * a = a * e = a$
            \item \textbf{وجود عضو وارون:}
            $\forall a \in G: \exists a^{-1} \in G: a * a^{-1} = a^{-1} * a = e$
        \end{enumerate}
        \label{def:group}
    \end{definition}
    
    \begin{example}
        مجموعه اعداد صحیح $\Z$ با عمل جمع معمولی یک گروه است:
        \begin{itemize}
            \item عضو خنثی: $0$
            \item وارون $a$: $-a$
        \end{itemize}
        \label{ex:z-group}
    \end{example}
    
    \begin{example}
        مجموعه اعداد حقیقی غیرصفر $\R \setminus \{0\}$
        با عمل ضرب معمولی یک گروه است:
        \begin{itemize}
            \item عضو خنثی: $1$
            \item وارون $a$: $\frac{1}{a}$
        \end{itemize}
        \label{ex:r-group}
    \end{example}
    
    \section{قضایای مقدماتی}
    
    \begin{theorem}[یکتایی عضو خنثی]
        عضو خنثی در هر گروه یکتاست.
        \label{thm:unique-identity}
    \end{theorem}
    
    \begin{proof}
        فرض کنید $e$ و $e'$ هر دو عضو خنثی گروه $G$ باشند.
        آنگاه:
        \[
        e = e * e' \quad \text{(چون $e'$ خنثی است)}
        \]
        \[
        e * e' = e' \quad \text{(چون $e$ خنثی است)}
        \]
        بنابراین $e = e'$ و عضو خنثی یکتاست.
    \end{proof}
    
    \begin{theorem}[یکتایی عضو وارون]
        وارون هر عضو در یک گروه یکتاست.
        \label{thm:unique-inverse}
    \end{theorem}
    
    \begin{proof}
        فرض کنید $a \in G$ و $b$ و $c$ دو وارون $a$ باشند:
        \begin{align*}
            b &= b * e \quad \text{(خنثی)} \\
            &= b * (a * c) \quad \text{($c$ وارون $a$ است)} \\
            &= (b * a) * c \quad \text{(شرکت‌پذیری)} \\
            &= e * c \quad \text{($b$ وارون $a$ است)} \\
            &= c \quad \text{(خنثی)}
        \end{align*}
        بنابراین $b = c$ و عضو وارون یکتاست.
    \end{proof}
    
    \section{تمرین‌ها}
    
    \begin{exercise}
        ثابت کنید که در هر گروه $G$، معادله $a * x = b$
        همواره دارای جواب یکتاست.
        \label{exc:group-equation}
    \end{exercise}
    
    \begin{solution}
        \textbf{وجود جواب:} $x = a^{-1} * b$ یک جواب است زیرا:
        \[
        a * (a^{-1} * b) = (a * a^{-1}) * b = e * b = b
        \]
        
        \textbf{یکتایی:} فرض کنید $x_{1}$ و $x_{2}$ دو جواب باشند:
        \begin{align*}
            a * x_{1} &= b \\
            a * x_{2} &= b
        \end{align*}
        با برابر قرار دادن و ضرب از چپ در $a^{-1}$:
        \[
        a^{-1} * (a * x_{1}) = a^{-1} * (a * x_{2})
        \implies x_{1} = x_{2}
        \]
        بنابراین جواب یکتاست.
    \end{solution}
    
    \begin{exercise}
        ثابت کنید که $(\Z_{n}, +_{n})$ یک گروه است
        (مجموعه اعداد صحیح به پیمانه $n$ با عمل جمع پیمانه‌ای).
        \label{exc:zn-group}
    \end{exercise}
    
    %----------------------------------------
    % Chapter 3: Cross-References Test
    %----------------------------------------
    \chapter{آزمون ارجاع متقابل}
    
    \section{ارجاع به تمام محیط‌ها}
    
    در این بخش، توانایی ارجاع به انواع مختلف محیط‌های
    علمی را آزمایش می‌کنیم:
    
    \begin{itemize}
        \item \textbf{قضیه:} \thmref{thm:fundamental} می‌گوید
        هر عدد طبیعی بزرگتر از ۱ به صورت حاصلضرب اعداد اول
        قابل تجزیه است.
        
        \item \textbf{تعریف:} طبق \defref{def:prime}، عدد اول
        عددی است که دقیقاً دو مقسوم‌علیه مثبت داشته باشد.
        
        \item \textbf{لم:} \lemref{lem:euclid} یک لم کلیدی
        در نظریه اعداد است.
        
        \item \textbf{نتیجه:} \corref{cor:euclid-general}
        تعمیم لم اقلیدس به چند عامل است.
        
        \item \textbf{گزاره:} \ref{prop:smallest-divisor}
        کوچکترین مقسوم‌علیه را معرفی می‌کند.
        
        \item \textbf{مثال:} \exref{ex:primes-under-50}
        اعداد اول کوچکتر از ۵۰ را فهرست کرده است.
        
        \item \textbf{نکته:} \ref{rem:euclid-importance}
        اهمیت لم اقلیدس را توضیح می‌دهد.
        
        \item \textbf{تمرین:} \exref{exc:group-equation}
        حل معادله در گروه را می‌طلبد.
    \end{itemize}
    
    \section{ارجاع بین فصلی}
    
    در \chref{chap:analysis} (فصل بعدی)،
    مفاهیم حد و پیوستگی را بررسی خواهیم کرد.
    
    همچنین در \exref{exc:zn-group} (فصل قبل)،
    از شما خواسته شد گروه $\Z_{n}$ را بررسی کنید.
    
    %----------------------------------------
    % Chapter 4: Real Analysis
    %----------------------------------------
    \chapter{مبانی آنالیز حقیقی}
    \label{chap:analysis}
    
    \section{حد و همگرایی}
    
    \begin{definition}[حد دنباله]
        دنباله $(a_{n})_{n=1}^{\infty}$ به عدد حقیقی $L$
        \textbf{همگرا} است اگر:
        \[
        \forall \varepsilon > 0, \;
        \exists N \in \N, \;
        \forall n \geq N: \abs{a_{n} - L} < \varepsilon
        \]
        در این صورت می‌نویسیم $\lim_{n \to \infty} a_{n} = L$.
        \label{def:limit}
    \end{definition}
    
    \begin{example}
        دنباله $a_{n} = \frac{1}{n}$ به صفر همگراست، زیرا:
        \[
        \forall \varepsilon > 0, \;
        \text{با انتخاب } N > \frac{1}{\varepsilon}, \;
        \forall n \geq N: \abs{\frac{1}{n} - 0} < \varepsilon
        \]
        \label{ex:convergent-1over}
    \end{example}
    
    \begin{theorem}[یکتایی حد]
        اگر دنباله‌ای همگرا باشد، حد آن یکتاست.
        \label{thm:unique-limit}
    \end{theorem}
    
    \begin{proof}
        فرض کنید $\lim a_{n} = L_{1}$ و $\lim a_{n} = L_{2}$
        با $L_{1} \neq L_{2}$. قرار دهید
        $\varepsilon = \frac{\abs{L_{1} - L_{2}}}{2} > 0$.
        
        طبق تعریف حد (\defref{def:limit})، عدد $N_{1}$ وجود دارد
        که برای $n \geq N_{1}$، $\abs{a_{n} - L_{1}} < \varepsilon$.
        همچنین عدد $N_{2}$ وجود دارد که برای $n \geq N_{2}$،
        $\abs{a_{n} - L_{2}} < \varepsilon$.
        
        برای $n \geq \max(N_{1}, N_{2})$:
        \begin{align*}
            \abs{L_{1} - L_{2}}
            &= \abs{(L_{1} - a_{n}) + (a_{n} - L_{2})} \\
            &\leq \abs{L_{1} - a_{n}} + \abs{a_{n} - L_{2}} \\
            &< \varepsilon + \varepsilon = \abs{L_{1} - L_{2}}
        \end{align*}
        که تناقض است. بنابراین $L_{1} = L_{2}$.
    \end{proof}
    
    \begin{remark}
        عکس قضیه \thmref{thm:unique-limit} لزوماً برقرار نیست:
        یکتایی حد شرط کافی برای همگرایی نیست.
        دنباله $a_{n} = (-1)^{n}$ حد ندارد،
        با اینکه اگر داشت یکتا می‌بود.
        \label{rem:limit-converse}
    \end{remark}
    
    %----------------------------------------
    % Summary
    %----------------------------------------
    \chapter{خلاصه نتایج}
    
    \section{آمار محیط‌های استفاده شده}
    
    \begin{table}[h]
        \centering
        \caption{محیط‌های علمی استفاده شده در این آزمون}
        \label{tab:environments-used}
        \begin{tabular}{|c|C{4cm}|c|}
            \hline
            \textbf{محیط} & \textbf{تعداد} & \textbf{وضعیت} \\
            \hline
            \lr{theorem} & قضیه اساسی حساب، یکتایی خنثی، یکتایی وارون، یکتایی حد & \textcolor{matingreen}{✅} \\
            \lr{definition} & عدد اول، گروه، حد دنباله & \textcolor{matingreen}{✅} \\
            \lr{lemma} & لم اقلیدس & \textcolor{matingreen}{✅} \\
            \lr{corollary} & تعمیم لم اقلیدس & \textcolor{matingreen}{✅} \\
            \lr{proposition} & کوچکترین مقسوم‌علیه & \textcolor{matingreen}{✅} \\
            \lr{example} & اعداد اول، ۶، $\Z$، $\R$، همگرایی & \textcolor{matingreen}{✅} \\
            \lr{remark} & اهمیت لم اقلیدس، ویژگی عدد ۲، عکس یکتایی حد & \textcolor{matingreen}{✅} \\
            \lr{proof} & قضیه اساسی، لم اقلیدس، یکتایی خنثی، یکتایی وارون، یکتایی حد & \textcolor{matingreen}{✅} \\
            \lr{exercise} & حل معادله در گروه، گروه $\Z_{n}$ & \textcolor{matingreen}{✅} \\
            \lr{solution} & حل معادله گروه & \textcolor{matingreen}{✅} \\
            \hline
        \end{tabular}
    \end{table}
    
    \section{نتیجه}
    
    \textbf{تمام محیط‌های علمی \lr{MatinBook} نسخه ۱ با موفقیت آزمایش شدند.}
    ۱۰ نوع جعبه رنگی، ارجاع‌های هوشمند، و اثبات‌های ریاضی
    همگی به درستی کار می‌کنند. 🎉
    
\end{document}